%-------------------------
% Resume in Latex
% Author : Pavel Kolyagin
% Based off of: https://github.com/sb2nov/resume
% License : MIT
%------------------------

\documentclass[letterpaper,11pt]{article}

\usepackage{latexsym}
\usepackage[empty]{fullpage}
\usepackage{titlesec}
\usepackage{marvosym}
\usepackage[usenames,dvipsnames]{color}
\usepackage{verbatim}
\usepackage{enumitem}
\usepackage[hidelinks]{hyperref}
\usepackage{fancyhdr}
\usepackage[T2A]{fontenc}
\usepackage[utf8]{inputenc}
\usepackage[english,russian]{babel}
\usepackage{tabularx}
\input{glyphtounicode}


%----------FONT OPTIONS----------
% sans-serif
% \usepackage[sfdefault]{FiraSans}
% \usepackage[sfdefault]{roboto}
% \usepackage[sfdefault]{noto-sans}
% \usepackage[default]{sourcesanspro}

% serif
% \usepackage{CormorantGaramond}
% \usepackage{charter}


\pagestyle{fancy}
\fancyhf{} % clear all header and footer fields
\fancyfoot{}
\renewcommand{\headrulewidth}{0pt}
\renewcommand{\footrulewidth}{0pt}

% Adjust margins
\addtolength{\oddsidemargin}{-0.5in}
\addtolength{\evensidemargin}{-0.5in}
\addtolength{\textwidth}{1in}
\addtolength{\topmargin}{-.5in}
\addtolength{\textheight}{1.0in}

\urlstyle{same}

\raggedbottom
\raggedright
\setlength{\tabcolsep}{0in}

% Sections formatting
\titleformat{\section}{
  \vspace{-4pt}\scshape\raggedright\large
}{}{0em}{}[\color{black}\titlerule \vspace{-5pt}]

% Ensure that generate pdf is machine readable/ATS parsable
\pdfgentounicode=1

%-------------------------
% Custom commands
\newcommand{\resumeItem}[1]{
  \item\small{
    {#1 \vspace{-2pt}}
  }
}

\newcommand{\resumeSubheading}[4]{
  \vspace{-2pt}\item
    \begin{tabular*}{0.97\textwidth}[t]{l@{\extracolsep{\fill}}r}
      \textbf{#1} & #2 \\
      \textit{\small#3} & \textit{\small #4} \\
    \end{tabular*}\vspace{-7pt}
}

\newcommand{\resumeSubSubheading}[2]{
    \item
    \begin{tabular*}{0.97\textwidth}{l@{\extracolsep{\fill}}r}
      \textit{\small#1} & \textit{\small #2} \\
    \end{tabular*}\vspace{-7pt}
}

\newcommand{\resumeProjectHeading}[2]{
    \item
    \begin{tabular*}{0.97\textwidth}{l@{\extracolsep{\fill}}r}
      \small#1 & #2 \\
    \end{tabular*}\vspace{-7pt}
}

\newcommand{\resumeSubItem}[1]{\resumeItem{#1}\vspace{-4pt}}

\renewcommand\labelitemii{$\vcenter{\hbox{\tiny$\bullet$}}$}

\newcommand{\resumeSubHeadingListStart}{\begin{itemize}[leftmargin=0.15in, label={}]}
\newcommand{\resumeSubHeadingListEnd}{\end{itemize}}
\newcommand{\resumeItemListStart}{\begin{itemize}}
\newcommand{\resumeItemListEnd}{\end{itemize}\vspace{-5pt}}

%-------------------------------------------
%%%%%%  RESUME STARTS HERE  %%%%%%%%%%%%%%%%%%%%%%%%%%%%


\begin{document}

%----------HEADING----------
\begin{center}
    \textbf{\Huge \scshape Колягин Павел} \\ \vspace{1pt}
    \small Software Engineer --- Internal Tools \& AI $|$ Опыт работы: 3 года 11 месяцев \\ \vspace{1pt}
    \small +7 914-978-82-75 $|$ \href{mailto:inbox@paulrainy.ru}{\underline{inbox@paulrainy.ru}} $|$
    \href{https://github.com/paulrainy}{\underline{github.com/paulrainy}}
\end{center}


%-----------ОБО МНЕ-----------
\section{Обо мне}
  \small{Инженер полного цикла: проектирую, разрабатываю и внедряю внутренние инструменты --- от бэкенда и БД до фронтенда и деплоя. Строю системы, которые ускоряют работу инженерных команд: анализ результатов нагрузочного тестирования, мониторинг тестовых стендов, генерация тестовых данных и автотестов. Применяю LLM и векторный поиск (RAG, эмбеддинги) в продуктовых задачах. Бэкграунд --- автоматизация тестирования (Java/Python) и разработка для защищённых систем (QNX, ФСТЭК).

  Формат работы: полная занятость, гибкий график, удаленная работа}


%-----------ОПЫТ РАБОТЫ-----------
\section{Опыт работы}
  \resumeSubHeadingListStart

    \resumeSubheading
      {Инженер по автоматизации тестирования}{Июнь 2024 --- настоящее время}
      {ООО «РСХБ-Автоматизация» и ООО «РСХБ-ИНТЕХ»}{Россия, Москва}
      \resumeItemListStart
        \resumeItem{Разработал fullstack-инструмент анализа результатов нагрузочного тестирования (Python + Streamlit + PostgreSQL): интеграция с 7 источниками метрик (Zabbix, Prometheus, InfluxDB, Grafana, Kibana, Oracle, PostgreSQL), автоматическое сравнение релизов с подсветкой отклонений, выгрузка отчетов в web, docx и xlsx}
        \resumeItem{Участвовал в разработке решения для hr отдела, позволяющее генерировать презентации и промо материалы с помощью больших языковых моделей. Добавил в проект поддержку Nano Banana и gpt-5-image.}
        \resumeItem{Разработал Down Detector как готовое fullstack решение (Python + FastAPI + TypeScript + React) для Rest сервисов всех тестовых окружений компании, повысив прозрачность состояния тестовых стендов для тестировщиков.}
        \resumeItem{Построил RAG-пайплайн для генерации автотестов на основе тест-кейсов (Python + Sentence-Transformers + Apache Airflow), в разы ускорив процесс автоматизации тест-кейсов автоматизаторами.}
        \resumeItem{Спроектировал 86 API тест-кейсов и автоматизировал 164 API и UI тест-кейса за второе полугодие 2025}
        \resumeItem{Исследовал возможности embedding и reranker-моделей в области актуализации пар тест-кейс --- автотест}
        \resumeItem{Выполнил рефакторинг legacy бэкенда сервиса по быстрой генерации клиента банка в тестовом окружении: разбил монолит на микросервисы, что улучшило читаемость и добавило возможность масштабирования функционала проекта}
        \resumeItem{Написал парсер для legacy почтового сервиса Lotus Notes (Java 21 + Selenide), облегчив пользователям миграцию на новый почтовый сервис, а также сохранив все их вложения}
      \resumeItemListEnd

    \resumeSubheading
      {Инженер}{Август 2022 --- Июнь 2024}
      {АО «Концерн „Моринсис-Агат``»}{Россия, Москва}
      \resumeItemListStart
        \resumeItem{Интегрировал СДЗ «Аккорд-АМДЗ» в информационную систему под управлением ОС QNX Neutrino}
        \resumeItem{Разработал пользовательский API для поддержки СДЗ «Аккорд-АМДЗ» в ОС QNX Neutrino}
        \resumeItem{Реализовал десктопное приложение для перевода геодезических координат между системами ПЗ90.11, СК-42 и WGS-84}
      \resumeItemListEnd

  \resumeSubHeadingListEnd


%-----------НАГРАДЫ И ВЫСТУПЛЕНИЯ-----------
\section{Награды и выступления}
  \resumeSubHeadingListStart
    \resumeSubheading
      {Победитель}{05.2026}
      {Конкурс «Лучший по профессии», номинация «Лучший тестировщик» (РСХБ-Автоматизация)}{Россия, Москва}
    \resumeSubheading
      {Спикер}{12.2025}
      {Testify \#12}{Россия, Москва}
    \resumeSubheading
      {Стендист}{04.2025}
      {Heisenbug Spring 2025}{Россия, Москва}
  \resumeSubHeadingListEnd


%-----------ОБРАЗОВАНИЕ-----------
\section{Образование}
  \resumeSubHeadingListStart
    \resumeSubheading
      {РТУ МИРЭА (Российский технологический университет «МИРЭА»)}{Россия, Москва}
      {Магистр, Информатика и вычислительная техника}{Сентябрь 2023 --- Июль 2025}
    \resumeSubheading
      {НИУ МЭИ (Национальный исследовательский университет «МЭИ»)}{Россия, Москва}
      {Бакалавр, Информационная безопасность}{Сентябрь 2019 --- Июль 2023}
  \resumeSubHeadingListEnd


%-----------ДОПОЛНИТЕЛЬНОЕ ОБРАЗОВАНИЕ-----------
\section{Дополнительное образование}
  \resumeSubHeadingListStart
    \resumeSubheading
      {Академия трансформации (АО «Россельхозбанк»)}{Россия, Москва}
      {Слушатель}{Апрель 2026}
      \resumeItemListStart
        \resumeItem{«Оптимизация процессов. Подход DMAIC» --- факультет «Управление процессами»}
      \resumeItemListEnd
    \resumeSubheading
      {РТУ МИРЭА (Российский технологический университет «МИРЭА»)}{Россия, Москва}
      {Студент}{Сентябрь 2023 --- Декабрь 2023}
      \resumeItemListStart
        \resumeItem{Менеджмент (профессиональная переподготовка)}
      \resumeItemListEnd
    \resumeSubheading
      {НИУ МЭИ (Национальный исследовательский университет «МЭИ»)}{Россия, Москва}
      {Студент}{Сентябрь 2022 --- Декабрь 2022}
      \resumeItemListStart
        \resumeItem{Цифровизация и развитие информационных систем (профессиональная переподготовка): разработка программного обеспечения на C++, блокчейн и криптовалюты, администрирование Linux и Windows (AD)}
      \resumeItemListEnd
  \resumeSubHeadingListEnd


%-----------СТЕК ТЕХНОЛОГИЙ-----------
\section{Стек технологий}
 \begin{itemize}[leftmargin=0.15in, label={}]
    \small{\item{
     \textbf{ML \& AI}{: Camel-ai, Sentence-transformers, Transformers, Scikit-learn, PyTorch} \\
     \textbf{Backend}{: Spring, FastAPI, Uvicorn, SQLAlchemy} \\
     \textbf{DevOps}{: Docker, Docker-compose, Apache Airflow, Github Actions} \\
     \textbf{Frontend}{: Vue3, Streamlit} \\
     \textbf{Programming}{: Java, Python, C++, C} \\
     \textbf{Databases}{: Milvus, QDrant, PostgreSQL, ELK Stack} \\
     \textbf{Языки}{: Русский (родной), Английский (B1)}
    }}
 \end{itemize}


%-------------------------------------------
\end{document}
